%% =====================================================================
%%  Corrigés — Unit 14 « What's the weather like? »
%% =====================================================================
\corrigesUnite{14}{What's the weather like?}

\begin{corrige}{1}
\begin{reponses}[3]
\item \rep{It's sunny.}
\item \rep{It's raining.}
\item \rep{It's snowing.}
\item \rep{It's cloudy.}
\item \rep{It's windy.}
\item \rep{It's foggy.}
\end{reponses}
\retiens{pour la pluie et la neige, on emploie le présent en \en{-ing} (\en{It's raining}) ; pour le
reste, un adjectif en \en{-y} (\en{sunny}, \en{cloudy}, \en{windy}, \en{foggy}).}
\end{corrige}

\begin{corrige}{2}
\begin{reponses}[4]
\item \rep{first}
\item \rep{second}
\item \rep{third}
\item \rep{fifth}
\item \rep{ninth}
\item \rep{twelfth}
\item \rep{twentieth}
\item \rep{thirty-first}
\end{reponses}
\piege{\en{nine} perd son \en{e} (\en{ninth}) ; \en{five} et \en{twelve} changent leur \en{ve} en \en{f}
(\en{fifth}, \en{twelfth}).}
\end{corrige}

\begin{corrige}{3}
\begin{reponses}[4]
\item \rep{on}
\item \rep{in}
\item \rep{at}
\item \rep{on}
\item \rep{in}
\item \rep{at}
\item \rep{in}
\item \rep{on}
\end{reponses}
\pourquoi{\en{at} + heure ; \en{on} + jour ou date ; \en{in} + mois, saison, partie de la journée. Une
exception à retenir : \en{at night}.}
\end{corrige}

\begin{corrige}{4}
\begin{reponses}[3][french]
\item \rep{V}
\item \rep{F}
\item \rep{V}
\item \rep{V}
\item \rep{F}
\item \rep{F}
\end{reponses}
\pourquoi{lundi, il pleut (\en{It's raining}) ; mardi, il fait 48\,°F, soit 9\,°C : c'est frais, pas
chaud ; vendredi, il neige et il fait −1\,°C : il gèle.}
\end{corrige}

\begin{corrige}{5}
\begin{reponses}[2]
\item \rep{November twenty-second}
\item \rep{August first}
\item \rep{June third}
\item \rep{September twelfth}
\item \rep{March fifteenth}
\item \rep{April thirtieth}
\end{reponses}
\piege{\en{6/3} n'est pas le 6 mars : aux États-Unis, le mois vient en premier. C'est le 3 juin.}
\end{corrige}

\begin{corrige}{6}
\begin{reponses}[6][french]
\item \rep{c}
\item \rep{e}
\item \rep{a}
\item \rep{f}
\item \rep{b}
\item \rep{d}
\end{reponses}
\retiens{\en{Thanksgiving} n'a pas de date fixe : c'est toujours le quatrième jeudi de novembre.}
\end{corrige}

\begin{corrige}{7}
\begin{reponses}[3]
\item \rep{short}
\item \rep{bird}
\item \rep{born}
\item \rep{turn}
\item \rep{word}
\item \rep{form}
\end{reponses}
\pourquoi{avec \phon{ɔːr}, les lèvres s'arrondissent nettement ; avec \phon{ɝː}, elles restent
détendues. Texte de la piste page~\pageref{trans-14.10}.}
\end{corrige}

\begin{corrige}{8}
\begin{reponses}[4]
\item \rep{cloudy}
\item \rep{50}
\item \rep{raining}
\item \rep{45}
\item \rep{windy}
\item \rep{38}
\item \rep{sunny}
\item \rep{82}
\end{reponses}
\pourquoi{texte de la piste page~\pageref{trans-14.11}. 82\,°F font environ 28\,°C : il fait chaud à
Miami, même en novembre !}
\end{corrige}

\begin{corrige}{9}
\begin{reponses}[4]
\item \rep{when's}
\item \rep{March}
\item \rep{yours}
\item \rep{summer}
\item \rep{June}
\item \rep{weather}
\item \rep{hot}
\item \rep{warm}
\end{reponses}
\retiens{\en{And yours?} = « Et le tien ? » : une façon simple de renvoyer la question.}
\end{corrige}

\begin{corrige}{10}
\begin{reponses}[1]
\item \rep{What's the weather like? — It's raining and it's cold.}
\item \rep{When's your birthday? — It's on June 3rd.}
\item \rep{It snows a lot in winter.}
\item \rep{It's thirty degrees Celsius: it's very hot!}
\item \rep{The concert is on Saturday, at eight o'clock.}
\item \rep{I don't like the wind.}
\item \rep{Christmas is on December 25th.}
\item \rep{I work in the morning.}
\end{reponses}
\piege{phrase 3 : une généralité, donc le présent simple (\en{It snows}) et non \en{It's snowing}.
Phrase 4 : jamais \en{It makes hot} !}
\end{corrige}

\begin{corrige}{11}
Réponse modèle :
\begin{anglais}
\begin{tcolorbox}[otBase,colback=otGreenTint,colframe=otGreenTint,fontupper=\fontsize{10.4}{14}\selectfont]
My favorite season is summer: June, July and August. In summer, it's sunny and hot. I go to the
beach and I swim every day. I love the long evenings! My birthday is in summer too: it's on July
12th.
\end{tcolorbox}
\end{anglais}
\retiens{\en{the beach} = « la plage » ; \en{the long evenings} = « les longues soirées ».}
\end{corrige}

\begin{corrige}{12}
Réponses modèles :
\begin{reponses}[1]
\item \rep{It's sunny and warm.} \textit{ou} \rep{It's cloudy and cold.}
\item \rep{Yes, it is. It's freezing!} \textit{ou} \rep{No, it isn't. It's cool.}
\item \rep{My favorite season is spring.}
\item \rep{It's on October 9th.}
\item \rep{It's November 18th.}
\end{reponses}
\end{corrige}

\begin{corrigetest}
\begin{reponses}[5][french]
\item \rep{c}
\item \rep{a}
\item \rep{b}
\item \rep{b}
\item \rep{a}
\item \rep{c}
\item \rep{a}
\item \rep{c}
\item \rep{b}
\item \rep{a}
\end{reponses}
\pourquoi{question 9 : (68 − 30) ÷ 2 ≈ 19, soit environ 20\,°C. Question 3 : en anglais, on ne « fait »
pas le temps : \en{It's cold}.}
\end{corrigetest}
